\documentclass[../../main.tex]{subfiles}

\begin{document}

\section{域、正性以及完备性公理}\label{section:ch01-域正性完备性}

假设给定实数集 $\mathbf{R}$, 使得对于每对实数 $a$ 和 $b$, 存在有定义的实数 $a+b$ 和 $ab$, 分别称为 $a$ 和 $b$ 的和与积. 它们满足以下的域公理、正性公理与完备性公理.

\begin{definition}[域公理]\label{definition:ch01-域公理}
\begin{enumerate}[(1)]
\item 加法的交换性: 对所有实数 $a$ 和 $b$,
\begin{align*}
a+b=b+a.
\end{align*}
\item 加法的结合性: 对所有实数 $a,b$ 和 $c$,
\begin{align*}
(a+b)+c=a+(b+c).
\end{align*}
\item 加法的单位元: 存在实数, 记为 $0$, 使得对所有实数 $a$,
\begin{align*}
0+a=a+0=a.
\end{align*}
\item 加法的逆元: 对每个实数 $a$, 存在实数 $b$ 使得
\begin{align*}
a+b=0.
\end{align*}
\item 乘法的交换性: 对所有实数 $a$ 和 $b$,
\begin{align*}
ab=ba.
\end{align*}
\item 乘法的结合性: 对所有实数 $a,b$ 和 $c$,
\begin{align*}
(ab)c=a(bc).
\end{align*}
\item 乘法的单位元: 存在实数, 记为 $1$, 使得对所有实数 $a$,
\begin{align*}
1a=a1=a.
\end{align*}
\item 乘法的逆元: 对每个实数 $a\neq 0$, 存在实数 $b$ 使得
\begin{align*}
ab=1.
\end{align*}
\item 分配性: 对所有实数 $a,b$ 和 $c$,
\begin{align*}
a(b+c)=ab+ac.
\end{align*}
\item 非平凡性假设:
\begin{align*}
1\neq 0.
\end{align*}
\end{enumerate}
\end{definition}

满足上述公理的任何集合称为域. 从加法的交换性可以得出加法的单位元 $0$ 是唯一的, 从乘法的交换性得出乘法的单位元 $1$ 也是唯一的. 加法的逆元和乘法的逆元也是唯一的. 我们记 $a$ 的加法的逆为 $-a$, 且若 $a\neq 0$, 记它的乘法逆为 $a^{-1}$ 或 $1/a$. 若有一个域, 我们能实施所有初等代数的运算, 包括解线性方程组. 我们不加声明地使用这些公理的多种推论.\footnote{域公理的推论的系统发展可在 Garrett Birkhoff 和 Saunders MacLane 的经典书《A Survey of Modern Algebra》\cite{BM97} 的第 1 章找到.}

\begin{definition}[正性公理]\label{definition:ch01-正性公理}
在实数中存在自然的序的概念: 大于, 小于, 等等. 编辑这些性质的方便方法是具体给出整数集满足的公理. 存在称为正数的实数集, 记为 $\mathcal{P}$. 它有以下两个性质:
\begin{enumerate}[(P1)]
\item 若 $a$ 和 $b$ 是正的, 则 $ab$ 和 $a+b$ 也是正的.
\item 对于实数 $a$, 以下三种情况恰有一种成立:
\begin{align*}
a\text{ 是正的},\qquad -a\text{ 是正的},\qquad a=0.
\end{align*}
\end{enumerate}
\end{definition}

正性公理以自然的方式引出实数的序: 对于实数 $a$ 和 $b$, 定义 $a>b$ 意味着 $a-b$ 是正的, 而 $a\geqslant b$ 意味着 $a>b$ 或 $a=b$. 定义 $a<b$ 意味着 $b>a$, 而 $a\leqslant b$ 意味着 $b\geqslant a$.

用域公理和正性公理, 可以正式证明不等式的常见性质（见习题 2）. 给定实数 $a$ 和 $b$ 满足 $a<b$, 我们定义 $(a,b)=\{x\mid a<x<b\}$, 且说 $(a,b)$ 的点落在 $a$ 与 $b$ 之间. 我们称非空实数集 $I$ 为区间, 若对 $I$ 中任意两点, 所有落在这两点之间的点也属于 $I$. 当然, 集合 $(a,b)$ 是区间. 以下集合也是区间:
\begin{align}
[a,b]=\{x\mid a\leqslant x\leqslant b\};\quad [a,b)=\{x\mid a\leqslant x<b\};\quad (a,b]=\{x\mid a<x\leqslant b\}.\label{align:ch01-区间}
\end{align}

\begin{definition}[完备性公理]\label{definition:ch01-完备性公理}
非空实数集 $E$ 称为有上界, 若存在实数 $b$, 使得对所有 $x\in E$, $x\leqslant b$. 数 $b$ 称为 $E$ 的上界. 类似地, 定义集合有下界以及一个数为集合的下界. 有上界的集合未必有最大的成员. 但下一个公断言它有一个最小的上界:

令 $E$ 为有上界的非空实数集. 则在 $E$ 的上界的集合中有一个最小的上界.
\end{definition}

对于实数的非空集 $E$, $E$ 的最小上界的存在性由完备性公理来保证, 记为 l.u.b.\,$E$. $E$ 的最小上界通常称为 $E$ 的上确界且记为 $\sup E$. 从完备性公理得出每个有下界的非空实数集 $E$ 有最大下界, 记为 g.l.b.\,$E$. 它常称为 $E$ 的下确界且记为 $\inf E$. 一个非空实数集称为有界的, 若它既有下界又有上界.

\begin{proposition}[三角不等式]\label{proposition:ch01-三角不等式}
定义实数 $x$ 的绝对值 $|x|$ 为: 若 $x\geqslant 0$ 则等于 $x$, 若 $x<0$ 则等于 $-x$. 以下不等式（称为三角不等式）在数学分析中是基本的: 对任何实数对 $a$ 和 $b$,
\begin{align*}
|a+b|\leqslant |a|+|b|.
\end{align*}
\end{proposition}

\begin{definition}[扩充的实数]\label{definition:ch01-扩充实数}
引入符号 $\infty$ 和 $-\infty$ 并对所有实数 $x$ 写 $-\infty<x<\infty$ 是方便的. 我们称集合 $\mathbf{R}\cup\{\pm\infty\}$ 为扩充的实数. 若非空实数集 $E$ 没有上界, 我们定义它的上确界为 $\infty$. 定义空集的上确界为 $-\infty$ 也是方便的. 因此每个实数集有一个属于扩充的实数的上确界. 类似地, 可以推广下确界的概念使得每个非空的实数集有属于扩充实数的下确界. 我们将定义实数序列的极限, 而允许极限是扩充的实数是方便的. 收敛到实数的实数序列的许多性质在极限是 $\pm\infty$ 时继续成立, 例如, 和的极限是极限的和且积的极限是极限的积, 若我们做出以下和与积的意义的推广: $\infty+\infty=\infty$, $-\infty-\infty=-\infty$, 对每个实数 $x$, $x+\infty=\infty$ 而 $x-\infty=-\infty$; 若 $x>0$, $x\cdot\infty=\infty$ 而 $x\cdot(-\infty)=-\infty$; 若 $x<0$, $x\cdot\infty=-\infty$ 而 $x\cdot(-\infty)=\infty$. 定义 $(-\infty,\infty)=\mathbf{R}$. 对于 $a,b\in\mathbf{R}$, 定义
\begin{align*}
(a,\infty)=\{x\in\mathbf{R}\mid a<x\},\qquad (-\infty,b)=\{x\in\mathbf{R}\mid x<b\}
\end{align*}
与
\begin{align*}
[a,\infty)=\{x\in\mathbf{R}\mid a\leqslant x\},\qquad (-\infty,b]=\{x\in\mathbf{R}\mid x\leqslant b\}.
\end{align*}
上面形式的集合是无界区间. 从 $\mathbf{R}$ 的完备性可以推出所有无界区间是上述形式的一种, 而所有有界区间是\eqref{align:ch01-区间}列出的形式, 带有形如 $(a,b)$ 的区间, 我们将其留作练习.
\end{definition}

\begin{exercise}\label{exercise:ch01-01}
对 $a\neq 0$ 和 $b\neq 0$, 证明 $(ab)^{-1}=a^{-1}b^{-1}$.
\end{exercise}

\begin{exercise}\label{exercise:ch01-02}
验证:
\begin{enumerate}[(i)]
\item 对每个实数 $a\neq 0$, $a^{2}>0$. 特别是, $1>0$, 由于 $1\neq 0$ 且 $1=1^{2}$.
\item 对每个正数 $a$, 它的乘法逆 $a^{-1}$ 也是正的.
\item 设 $a>b$, 有
\begin{align*}
\text{若 }c>0,\quad ac>bc;\qquad\text{而若 }c<0,\quad ac<bc.
\end{align*}
\end{enumerate}
\end{exercise}

\begin{exercise}\label{exercise:ch01-03}
对于非空实数集 $E$, 证明 $\inf E=\sup E$ 当且仅当 $E$ 由单点组成.
\end{exercise}

\begin{exercise}\label{exercise:ch01-04}
令 $a$ 和 $b$ 为实数.
\begin{enumerate}[(i)]
\item 证明: 若 $ab=0$, 则 $a=0$ 或 $b=0$.
\item 验证 $a^{2}-b^{2}=(a-b)(a+b)$, 并从 (i) 部分推出: 若 $a^{2}=b^{2}$, 则 $a=b$ 或 $a=-b$.
\item 令 $c$ 为正实数. 定义 $E=\{x\in\mathbf{R}\mid x^{2}<c\}$. 验证 $E$ 是非空的且有上界. 定义 $x_{0}=\sup E$. 证明 $x_{0}^{2}=c$. 用 (ii) 部分证明存在唯一的 $x>0$ 使得 $x^{2}=c$. 记之为 $\sqrt{c}$.
\end{enumerate}
\end{exercise}

\begin{exercise}\label{exercise:ch01-05}
令 $a,b$ 和 $c$ 为实数, 使得 $a\neq 0$. 考虑二次方程
\begin{align*}
ax^{2}+bx+c=0,\quad x\in\mathbf{R}.
\end{align*}
\begin{enumerate}[(i)]
\item 假定 $b^{2}-4ac>0$. 用域公理和前一个习题来配方, 从而证明该方程恰有由
\begin{align*}
x=\frac{-b+\sqrt{b^{2}-4ac}}{2a}\quad\text{和}\quad x=\frac{-b-\sqrt{b^{2}-4ac}}{2a}
\end{align*}
给出的两个解.
\item 现在假定 $b^{2}-4ac<0$. 证明二次方程无解.
\end{enumerate}
\end{exercise}

\begin{exercise}\label{exercise:ch01-06}
用完备性公理证明每个有下界的非空实数集有下确界, 且
\begin{align*}
\inf E=-\sup\{-x\mid x\in E\}.
\end{align*}
\end{exercise}

\begin{exercise}\label{exercise:ch01-07}
对实数 $a$ 和 $b$, 验证:
\begin{enumerate}[(i)]
\item $|ab|=|a|\,|b|$.
\item $|a+b|\leqslant |a|+|b|$.
\item 对 $\varepsilon>0$,
\begin{align*}
|x-a|<\varepsilon\text{ 当且仅当 }a-\varepsilon<x<a+\varepsilon.
\end{align*}
\end{enumerate}
\end{exercise}

\end{document}